%%%PAGE 189 rot=0 legibility=good summary=易错点清单（第二页，序号1-9）：核能单位换算、卫星轨道Ⅰ Ⅱ Ⅲ比较、行星绕火星求半径、核反应与裂变聚变、变轨、物质波与光伏效应、同步卫星高度、激光测火星半径、能级与光电效应、测火星密度
% 说明：页首右上角的淡铅笔化学字迹（n CO2 … 燃烧得CO和CO2混合气体 ρ=1.429 g·L^-1；CO为0.4×3/4×22.4=6.72 L）与右上角翘起的另一页（No./Date、“…60 g/mL的容…”）在 p186 的同一位置也出现，应是垫在笔记本下面的另一张纸/页的透印，不是本页内容，未作为本页内容转录。
%%%BLOCK id=189-01 ch=17 sec=核能·质量亏损与单位换算 kind=note alt=none cont=none y=0.08-0.16
1、1个\unclear{核反应}释放\unclear{的} $E_0$ 能量\quad $\mathrm{MeV}=E_0\times1.6\times10^{-13}\unit{J}$
%%%END
%%%BLOCK id=189-02 ch=05 sec=卫星变轨·椭圆轨道与圆轨道比较 kind=note alt=none cont=none y=0.11-0.31
% 本块的圆/椭圆图与上一行 1、 的文字在页面上有重叠（图中的大圆为较淡的“构造轨道”）；图中有两条细长斜线自页顶贯穿到图内
2、

%FIG p189_a 0.0 0.105 0.46 0.31 soft
\notefig[0.5]{figures/p189_a.png}

$S_2<S_3$（$T$ 同）

$\mathrm{II}$ $\mathrm{III}$ 在 $P$ 点 $a$ 相同 但 $a_{\text{向}}$ 不同 ☆% CHECK: “a向”二字写得潦草，疑为 a_向（向心加速度）；原稿“a向不同”下有下划线

$V_{2E}<V_4<V_3<V_1<V_{2B}$
%%%END
%%%BLOCK id=189-03 ch=05 sec=天体质量·绕火星运动半径 kind=problem alt=none cont=none y=0.31-0.47
\begin{problem}[3]
行星绕火星以 $v$ 匀圆周运动．$M_{\text{地}}=QM_{\text{火}}$．$R_{\text{地}}=R$．$g_{\text{地}}=g$．求绕火星运动半径？
\begin{solution}
$\dfrac{GMm}{R^2}=mg$\quad $GM=gR^2$

$M_{\text{火}}=\dfrac{M_{\text{地}}}{Q}$

$\dfrac{G\dfrac{M_{\text{地}}}{Q}m_1}{r^2}=m_1\dfrac{v^2}{r}$\quad $r=\dfrac{GM}{Qv^2}=\dfrac{gR^2}{Qv^2}$
\end{solution}
\end{problem}
%%%END
%%%BLOCK id=189-04 ch=17 sec=核反应·裂变与聚变 kind=note alt=none cont=none y=0.44-0.52
% 原稿对“热核反应”“核聚变”“除氢弹外全为核裂变”“镉棒控制核反应速度”“临界质量才裂变”有下划线，“铀块”与“大于”各被圈出；“放能”被红叉划去
4、热核反应是核聚变．除氢弹外全为核裂变\quad 镉棒控制核反应速度．（\unclear{铀块}大于临界质量才裂变）

一切核反应都\struck{放能}\quad \red{只有有质量亏损的反应才放能}
%%%END
%%%BLOCK id=189-05 ch=05 sec=卫星变轨·高轨与低轨 kind=note alt=none cont=none y=0.50-0.54
5、高轨 $\to$ 低轨 要减速制动 向前喷气 $E_{\text{机}}\downarrow$% 原稿“减速制动 向前喷气 E机↓”下有下划线
%%%END
%%%BLOCK id=189-06 ch=17 sec=物质波·动能与动量关系 kind=formula alt=none cont=none y=0.53-0.58
$E=\dfrac{p^2}{2m}$，\quad $p=\sqrt{2mE}$，\quad $\lambda=\dfrac hp=\dfrac{h}{\sqrt{2mE}}$
%%%END
%%%BLOCK id=189-07 ch=17 sec=光电效应·光伏效应 kind=note alt=none cont=none y=0.55-0.61
硅吸收光子 不是光电效应．PN 结作用下形成电流——光伏效应% 原稿“不是光电效应”及“PN结…光伏效应”下有下划线
%%%END
%%%BLOCK id=189-08 ch=05 sec=同步卫星·轨道半径与高度 kind=note alt=none cont=none y=0.57-0.66
% 序号 6 被页边遮挡，据顺序补为 6
6、$r=\sqrt[3]{\dfrac{GMT^2}{4\pi^2}}$\quad $\dfrac{r_{\text{火}}}{r_{\text{地}}}=\sqrt[3]{\dfrac{M_{\text{火}}T_{\text{火自}}^{2}}{M_{\text{地}}T_{\text{地自}}^{2}}}$\quad 但卫星高度 $h=r-R$．故卫星高度比不等于轨道半径之比 ☆% 原稿“h=r-R．故卫星高度比不等于轨道半径之比”下有下划线

$V=\sqrt{\dfrac{GM}{R}}$\quad\quad $g=\dfrac{GM}{R^2}$
%%%END
%%%BLOCK id=189-09 ch=05 sec=天体质量·探测器激光测距 kind=problem alt=none cont=none y=0.66-0.78
\begin{problem}[7]
探测器向火星表面发射一束激光，经过时间 $t$，收到传回的信号．相邻两次看到日出的时间间隔为 $T$．测得探测器速度为 $V$．（其下一行右侧注：绕火星转的周期为 $T$）% 原稿“相邻两次看到日出的时间间隔为T”下有下划线
\begin{solution}
（行间小字：探测器质量，）算不了行星质量（消去了），% CHECK: 小字“探测器质量，”写在下一行行首之上，语义疑为“探测器质量算不了，……（消去了）”

探测器到火星表面距离为 $h=\dfrac{ct}{2}$\quad $v=\dfrac{2\pi r}{T}$\quad $r=\dfrac{vT}{2\pi}$

则火星半径 $R=r-h=\dfrac{vT}{2\pi}-\dfrac{ct}{2}$\quad $\dfrac{GMm}{r^2}=\dfrac{mv^2}{r}$\quad $M=\dfrac{v^3T}{2\pi G}$% CHECK: 末式分子处有涂黑的墨团，T 上方似有横线
\end{solution}
\end{problem}
%%%END
%%%BLOCK id=189-10 ch=17 sec=光子能量与轨道能量 kind=note alt=05 cont=none y=0.78-0.81
% 原稿此块被划去：一条横线穿过“E=hc/λ=hν”并延伸到“低轨”二字，“E小”未被划
\struck{$E=\dfrac{hc}{\lambda}=h\nu$}\quad 低轨 $E$ 小
%%%END
%%%BLOCK id=189-11 ch=17 sec=氢原子能级·光电效应 kind=note alt=none cont=none y=0.81-0.86
% 序号被页边遮挡（疑为 8）
$n=2\to n=1$ 放出光子恰发生光电效应．从 $n=3\to n=1$\quad $E_{k\text{初}\max}=E_3-E_2$

$W_0=E_2-E_1$\quad\quad $E_k=E_3-E_1-(E_2-E_1)$
%%%END
%%%BLOCK id=189-12 ch=05 sec=天体密度·探测器绕火星 kind=problem alt=none cont=none y=0.86-1.00
\begin{problem}[9]
（行间小字：探测器绕火星）飞 $N$ 圈用时 $t$．
\begin{solution}
$T=\dfrac tN$\quad $V=\dfrac{2\pi r}{T}=\dfrac{2\pi Nr}{t}$\quad $a_{\text{向}}=\dfrac{v^2}{r}=\dfrac{4\pi^2N^2r}{t^2}$% CHECK: 末式分母 t^2 前似有一短横

$M_{\text{地}}=M$\quad $R_{\text{地}}=R$\quad $r_{\text{火}}=r$\quad $g_{\text{地}}=g$．

$\dfrac{GM_{\text{火}}m}{r^2}=\dfrac{mv^2}{r}$\quad $M_{\text{火}}=\rho\,\dfrac43\pi r^3$\quad $GM_{\text{火}}=gR^2$\quad $\therefore\rho_{\text{火星}}=\dfrac{3\pi MN^2}{gR^2t^2}$% CHECK: “GM_火=gR^2”中的下标似为“火”，按上文应为地球的 GM_地=gR^2

（小字：探测器质量约去了）

自转周期求不了\quad 求 $\rho$ 要有 $m$% CHECK: 末字可能是 M（地球质量）
\end{solution}
\end{problem}
%%%END
%%%PAGEEND 189
